\chapter{Tổng quan đề tài}

\section{Bối cảnh và Vấn đề}

Sự phát triển của các nền tảng học trực tuyến giúp giảng viên phân phối tài
liệu, bài tập và bài kiểm tra nhanh hơn so với lớp học truyền thống. Tuy
nhiên, trong nhiều học phần kỹ thuật, quy trình chuẩn bị nội dung vẫn phụ
thuộc nhiều vào thao tác thủ công: giảng viên đọc giáo trình, slide, đề cương
môn học, tóm tắt lại thành bài học ngắn, sau đó tự soạn câu hỏi kiểm tra và
đáp án. Khi tài liệu thay đổi hoặc môn học có nhiều chương, quy trình này tốn
thời gian, khó duy trì chất lượng đồng đều và khó cá nhân hóa cho từng nhóm
học viên.

Ở phía học viên, tài liệu dài gây khó khăn cho việc ôn tập theo từng đơn vị
kiến thức nhỏ. Học viên thường cần nội dung ngắn gọn, câu hỏi tự kiểm tra và
phản hồi tức thì để biết phần nào đã hiểu, phần nào cần học lại. Nhu cầu này
phù hợp với Micro-learning, trong đó nội dung học tập được chia thành các đơn
vị nhỏ, dễ tiếp cận, dễ lặp lại và có thể đánh giá thường xuyên.

Các mô hình ngôn ngữ lớn (Large Language Models --- LLMs) và Generative AI mở
ra khả năng tự động hóa một phần quy trình tạo nội dung học tập. LLM có thể
tóm tắt, diễn giải, sinh câu hỏi và giải thích đáp án. Tuy vậy, nếu gọi LLM
trực tiếp trên tài liệu thô, hệ thống dễ gặp các vấn đề như sinh nội dung
không bám nguồn, câu hỏi không gắn với chuẩn đầu ra học phần, thiếu kiểm soát
chất lượng và khó tích hợp vào Learning Management System (LMS) hiện hữu.

Từ đó, đề tài đặt ra bài toán: xây dựng một hệ thống tích hợp LMS có khả năng
chuyển đổi tài liệu học thuật thành Micro-Content và Quiz bằng Generative AI,
đồng thời vẫn đảm bảo khả năng kiểm duyệt học thuật, truy xuất nguồn tri thức
và gắn kết với chuẩn đầu ra học phần.

Các thách thức chính gồm:
\begin{itemize}
    \item \textbf{Tài liệu đầu vào không đồng nhất:} PDF, DOCX, PPTX có cấu
          trúc, heading, bảng biểu và định dạng khác nhau.
    \item \textbf{Ngữ cảnh dài:} LLM không nên nhận toàn bộ tài liệu một lần;
          hệ thống cần chunking và truy xuất phù hợp.
    \item \textbf{Độ tin cậy nội dung:} nội dung sinh ra phải bám sát tài
          liệu gốc và cần có bước Human-in-the-Loop để giảng viên review.
    \item \textbf{Liên kết chuẩn đầu ra:} Quiz cần đo đúng Learning Outcome
          thay vì chỉ hỏi rời rạc theo từng đoạn văn bản.
    \item \textbf{Tích hợp LMS:} hệ thống phải hoạt động bên cạnh Open edX
          hoặc Moodle, tránh buộc cơ sở đào tạo thay đổi toàn bộ LMS.
    \item \textbf{Đánh giá chất lượng:} cần đánh giá cả chức năng, truy xuất,
          chất lượng nội dung sinh ra, hiệu năng và trải nghiệm sử dụng.
\end{itemize}

\section{Mục tiêu đề tài}

\subsection{Mục tiêu tổng quát}

Xây dựng và đánh giá một hệ thống LMS tích hợp Generative AI có khả năng tự
động chuyển đổi tài liệu học thuật thành Micro-Content và Quiz, hỗ trợ review
của giảng viên, tìm kiếm ngữ nghĩa, truy xuất dựa trên RAG/GraphRAG và liên
kết nội dung với chuẩn đầu ra học phần.

\subsection{Mục tiêu cụ thể}

\begin{enumerate}
    \item Thiết kế kiến trúc hệ thống gồm frontend, backend gRPC, worker xử lý
          bất đồng bộ, PostgreSQL, Qdrant, Neo4j, MinIO và Redis/Queue.
    \item Xây dựng pipeline xử lý tài liệu gồm trích xuất nội dung, OCR khi
          cần, chunking, embedding, retrieval và generation.
    \item Thiết kế Curriculum-Aware Layer để trích xuất Course, Chapter,
          Learning Outcome và ánh xạ nội dung học tập với chuẩn đầu ra.
    \item Tích hợp hệ thống với LMS thông qua LTI 1.3/OIDC, ưu tiên Open edX
          và phân tích khả năng tương thích với Moodle.
    \item Xây dựng giao diện cho giảng viên review/publish nội dung và học
          viên học card/làm quiz.
    \item Đề xuất phương pháp đánh giá chức năng, tích hợp, chất lượng truy
          xuất, chất lượng nội dung sinh ra và hiệu năng hệ thống.
\end{enumerate}

\section{Phạm vi và Giới hạn}

Trong phạm vi luận văn, hệ thống tập trung vào các nội dung sau:
\begin{itemize}
    \item Tài liệu đầu vào dạng PDF, DOCX, PPTX và đề cương môn học dạng văn
          bản/PDF có cấu trúc tương đối rõ.
    \item Sinh Micro-Content, Flashcard/Lesson Card, câu hỏi trắc nghiệm và
          giải thích đáp án.
    \item Tìm kiếm ngữ nghĩa và RAG dựa trên chunk đã embedding.
    \item Knowledge Graph cho Course, Chapter, Learning Outcome, Assessment
          và quan hệ chunk hỗ trợ LO.
    \item Tích hợp LMS bằng LTI 1.3 với Open edX; Moodle được phân tích như
          hướng triển khai tương thích.
    \item Triển khai local/staging bằng Docker Compose.
\end{itemize}

Các giới hạn hiện tại gồm:
\begin{itemize}
    \item OCR production cho tài liệu scan chất lượng thấp chưa phải trọng tâm
          hiện thực; phần này được trình bày như thiết kế mở rộng.
    \item Kubernetes production đầy đủ chưa được triển khai trong phạm vi luận
          văn; hệ thống chỉ nêu định hướng chuyển đổi từ Docker Compose.
    \item GraphRAG nhiều bước và reasoning trên graph chưa được đánh giá như
          một module hoàn chỉnh.
    \item Đánh giá pilot quy mô lớn chưa thuộc phạm vi chính; các chỉ số chưa
          đo thực nghiệm được đánh dấu rõ để tránh kết luận vượt quá dữ liệu.
\end{itemize}

\section{Ý nghĩa thực tiễn}

Đề tài có ý nghĩa thực tiễn ở ba khía cạnh. Thứ nhất, hệ thống hỗ trợ giảng
viên giảm thời gian chuẩn bị học liệu ngắn và câu hỏi kiểm tra, đồng thời vẫn
giữ quyền kiểm duyệt trước khi công bố cho học viên. Thứ hai, học viên có thể
tiếp cận tài liệu theo từng đơn vị nhỏ, làm quiz và nhận phản hồi nhanh để ôn
tập hiệu quả hơn. Thứ ba, cách tiếp cận tích hợp qua LTI giúp hệ thống có thể
được đưa vào LMS hiện hữu như Open edX hoặc Moodle mà không phải thay thế hạ
tầng học tập đang sử dụng.

Về mặt kỹ thuật, đề tài minh họa cách kết hợp RAG, Knowledge Graph, Bloom
taxonomy và kiến trúc xử lý bất đồng bộ để xây dựng một pipeline AI có khả
năng truy xuất nguồn, kiểm soát chất lượng và mở rộng trong môi trường giáo
dục.

\section{Cấu trúc luận văn}

Luận văn được tổ chức thành 7 chương:
\begin{itemize}
    \item \textbf{Chương 1 --- Tổng quan đề tài:} trình bày bối cảnh, vấn đề,
          mục tiêu, phạm vi, giới hạn và ý nghĩa thực tiễn.
    \item \textbf{Chương 2 --- Nền tảng lý thuyết và công nghệ:} trình bày
          Micro-learning, Bloom taxonomy, LLM, Generative AI, RAG, GraphRAG,
          LTI 1.3 và các công nghệ lõi.
    \item \textbf{Chương 3 --- Phân tích yêu cầu và nghiệp vụ:} phân tích tác
          nhân, yêu cầu hệ thống, use case và luồng nghiệp vụ chính.
    \item \textbf{Chương 4 --- Thiết kế hệ thống:} trình bày kiến trúc tổng
          thể, tích hợp LMS, AI pipeline, thiết kế dữ liệu, Knowledge Graph,
          giao tiếp nội bộ và system flows.
    \item \textbf{Chương 5 --- Hiện thực và triển khai:} mô tả môi trường phát
          triển, frontend, backend, workers, tích hợp LMS và đóng gói triển
          khai.
    \item \textbf{Chương 6 --- Đánh giá hệ thống:} trình bày phương pháp đánh
          giá chức năng, tích hợp, retrieval, chất lượng nội dung sinh ra và
          hiệu năng.
    \item \textbf{Chương 7 --- Kết luận và hướng phát triển:} tổng kết kết
          quả, hạn chế, định hướng mở rộng và kế hoạch triển khai giai đoạn
          luận văn tốt nghiệp.
\end{itemize}
